\documentclass[11pt]{article}
\usepackage{amsmath,amssymb,amsfonts,amsthm}
\usepackage{geometry}
\usepackage{hyperref}
\usepackage{bm}
\geometry{margin=1in}

\title{Phase Mirror Multiplicity Agents (PMMA) \\
\large A Governance-First Architecture for Agentic Domain-Specific Reasoning}
\author{Citizen Gardens \\ The foundation of Multiplicity}
\date{April 29, 2026}

\begin{document}
\maketitle

\begin{abstract}
This defensive publication discloses the Phase Mirror Multiplicity Agents (PMMA) framework: a governance-first, mathematically grounded architecture for agentic domain-specific reasoning.
PMMA instantiates the tensions named in the ``Agentic Domain-Specific Reasoning'' Phase Mirror white paper---between autonomy and governance, probabilistic outputs and binary compliance, domain precision and platform generalization---as a recursively stabilized dynamical system with explicit causal validity gates, reliability budgets, and non-convergence escalation.
The framework is designed to be falsifiable via simulation and to act as a safety--efficiency manifold parameterized by enterprise policy levers rather than a single fixed operating point.
\end{abstract}
\newpage
\tableofcontents
\section*{Executive Summary}

Agentic AI systems promise autonomous decision-making but introduce structural tensions between autonomy, compliance, liability, and data quality that cannot be ``solved'' in a static way.
The Phase Mirror framework correctly identifies these tensions as permanent dissonances that must be governed through explicit bindings (specification of acceptable error, liability tiers, audit and kill-switch mechanisms) rather than informal alignment rhetoric.

Phase Mirror Multiplicity Agents (PMMA) operationalize this perspective by:
\begin{itemize}
  \item Representing each structural tension as a \emph{prime-indexed multiplicity component}, e.g.\ $p=2$ for autonomy vs.\ governance, $p=5$ for causal validity vs.\ no-code mis-specification risk.
  \item Evolving agent state via a recursive update map that encodes governance policies, reliability budgets, and causal validity gates.
  \item Enforcing \emph{admissible recursion}: autonomous recursion is allowed only when causal validity exceeds a threshold and reliability budget remains above a policy floor.
  \item Defining a violation functional and error budget $E$ that measure safety performance, and an intervention ratio that measures governance cost.
  \item Treating the architecture as a tunable safety--efficiency manifold: parameters can be calibrated to move from ``safe but inefficient'' regimes to enterprise-acceptable frontiers (e.g.\ $<5\%$ violations and $<25\%$ interventions).
\end{itemize}

The framework is designed to be implemented as:
\begin{enumerate}
  \item A simulation and test harness on a domain-specific pilot (e.g.\ lead qualification).
  \item A library of governance and causal validity operators with configuration-driven thresholds.
  \item A set of reproducible experiments comparing baseline human-in-the-loop, PMMA Phase I, and PMMA Phase II (tuned) policies under varying confounder omission and random seeds.
\end{enumerate}

This document serves both as a conceptual report and as a defensive publication, describing the architecture, mathematical formalism, and example code for PMMA, thereby placing the core ideas into the public domain as prior art.

\section{Conceptual Overview}

\subsection{From Named Tensions to Dynamical Systems}

The underlying white paper introduces three key families of tension:
\begin{enumerate}
  \item \textbf{Autonomy vs.\ Governance}: Agents need permission to fail, but enterprises demand deterministic guarantees, audit trails, and human accountability.
  \item \textbf{Probabilistic Outputs vs.\ Binary Compliance}: Bayesian or probabilistic reasoning conflicts with legal and compliance frameworks that encode binary categories (violation vs.\ non-violation, compliant vs.\ non-compliant).
  \item \textbf{Domain Precision vs.\ Platform Generalization}: Rich, causal, domain-specific models are in tension with scalable, general-purpose AI platforms, especially under no-code interfaces that hide theoretical complexity.
\end{enumerate}

The Phase Mirror methodology aims to replace vague claims with concrete \emph{mechanisms}: error budgets, liability tiers, audit pass rates, causal validity scores, violation rates, and kill-switches.
PMMA extends this by embedding those mechanisms into a stateful, recursive system, so that:
\begin{itemize}
  \item Each run of the agent modifies its governance posture (autonomy level, strictness, budget) in response to observed violations and causal validity.
  \item Policy parameters (thresholds, update magnitudes) act as controls that move the system along a safety--efficiency tradeoff curve.
  \item Non-convergence or severe causal degradation triggers escalation to human-on-exception, with explicit logged rationale.
\end{itemize}

\subsection{Prime-Indexed Multiplicity View}

To structurally organize the tensions, PMMA uses a prime-indexed decomposition.
Let $\Pi$ be a finite set of primes, and associate each
$p\in\Pi$ with a specific dissonance:
\begin{itemize}
  \item $p=2$: Autonomy vs.\ Governance.
  \item $p=5$: Causal Validity vs.\ No-Code Mis-Specification.
  \item Additional primes could encode other tensions (e.g.\ transparency vs.\ proprietary infrastructure, human-in-the-loop vs.\ human-on-exception).
\end{itemize}

The agent's state is then interpreted as living in a \emph{multiplicity space}
\[
  \mathcal{M}_t = \bigoplus_{p \in \Pi} \mu_p(t) \cdot \iota_p,
\]
where $\mu_p(t)$ is a multiplicity weight for prime $p$ at time $t$, and $\iota_p$ is an identity/interface operator describing the semantics of that prime (e.g.\ which variables and thresholds are active).
In practice, we will work with a more concrete state vector (Section~\ref{sec:math-overview}), but the multiplicity view clarifies that multiple tensions evolve together and may interact.

\section{Mathematical Overview}
\label{sec:math-overview}

\subsection{Core State and Policy Parameters}

At each discrete time step $t$, PMMA maintains a core state vector
\[
  M_t = \bigl(a_t, g_t, b_t, V_t, e_t\bigr),
\]
where:
\begin{itemize}
  \item $a_t \in [0,1]$: autonomy level (how willing is the system to act without human intervention).
  \item $g_t \in [0,1]$: governance strictness (how strongly constraints and review requirements are enforced).
  \item $b_t \in [0,1]$: reliability budget remaining (fraction of allowable failure budget not yet consumed).
  \item $V_t \in [0,1]$: causal validity score (quality of the current causal model or graph).
  \item $e_t \in \{\text{none},\text{human-review},\text{human-on-exception}\}$: exception routing state.
\end{itemize}

The policy configuration is given by:
\[
  \Theta = (\tau_V, V_\mathrm{critical}, \epsilon, \delta, k_\mathrm{max}, \text{decay rates, replenish rates, thresholds}),
\]
where:
\begin{itemize}
  \item $\tau_V$: admissible recursion threshold for causal validity (e.g.\ $0.6$ or $0.7$).
  \item $V_\mathrm{critical}$: severe degradation threshold (e.g.\ $0.5$).
  \item $\epsilon$: target error budget (e.g.\ $0.05$ for $5\%$).
  \item $\delta$: convergence tolerance for recursion (e.g.\ Euclidean distance).
  \item $k_\mathrm{max}$: maximum recursion steps before non-convergence escalation.
\end{itemize}

\subsection{Admissible Recursion and Causal Gate}

Given $M_t$, PMMA first determines whether autonomous recursion is allowed via a \emph{causal validity gate}.
Define:
\[
  \Lambda_t = \mathbf{1}[V_t \ge \tau_V].
\]
If $\Lambda_t = 0$, the system does not permit a fully autonomous update; instead it routes decisions towards human-review or human-on-exception.
This implements the paper's requirement that causal models be valid before deployment and that low-quality causal structure triggers additional oversight.

\subsection{Transition Dynamics}

We define a recursion map
\[
  M_{t+1} = \Pi(M_t; \Theta, \text{data}_t),
\]
where $\Pi$ depends on the current state, policy parameters, and observed data at time $t$ (including violations and validity gaps).
The transition decomposes into three components:

\paragraph{1. Autonomy and Governance Update.}

Let $v_t \in [0,1]$ be the empirical \emph{violation pressure} at time $t$ (fraction of decisions classified as violations in a recent batch), and let $u_t = \max(0,\tau_V - V_t)$ be a \emph{validity pressure} term.

We define soft, saturating updates:
\begin{align*}
  g_{t+1} &= \mathrm{clip}_{[0,1]}\bigl(g_t + \alpha_g (v_t, u_t)\, (1 - g_t)\bigr), \\
  a_{t+1} &= \mathrm{clip}_{[0,1]}\bigl(a_t - \alpha_a (v_t, u_t) + \beta_a(b_t)\bigr),
\end{align*}
where:
\begin{itemize}
  \item $\alpha_g$ increases governance strictness with violations and low validity, but is capped so that $g_t$ asymptotically approaches $1$ without large jumps.
  \item $\alpha_a$ decreases autonomy under violations and poor validity.
  \item $\beta_a(b_t)$ allows autonomy to increase slightly when budget is healthy (e.g.\ $b_t > 0.5$).
\end{itemize}

\paragraph{2. Reliability Budget Update.}

The reliability budget $b_t$ decays with violations and low validity, and replenishes slightly in ``clean'' periods:
\[
  b_{t+1} = \mathrm{clip}_{[0,1]}\bigl(b_t - \gamma_v v_t - \gamma_u u_t + r_t\bigr),
\]
where $\gamma_v,\gamma_u > 0$ are decay rates, and
\[
  r_t = \begin{cases}
    r_{\max} & \text{if } v_t \approx 0, \; u_t = 0, \\
    0 & \text{otherwise}.
  \end{cases}
\]
This implements an \emph{error budget} concept: the system may ``spend'' some allowable violation risk but must replenish by running clean before regaining autonomy.

\paragraph{3. Exception Routing and Non-Convergence.}

Within a recursion episode, we repeatedly apply $\Pi$ up to $k_\mathrm{max}$ steps or until
\[
  \|M_{t+1} - M_t\|_2 \le \delta.
\]
If convergence occurs, and $V_{t+1} \ge \tau_V$ and $b_{t+1} > 0$, the system may remain in an autonomous regime for the next batch.

If non-convergence occurs, or if $V_{t+1} < V_\mathrm{critical}$, PMMA routes to human-on-exception:
\[
  e_{t+1} = \begin{cases}
    \text{human-on-exception}, & \text{if non-convergence or } V_{t+1} < V_\mathrm{critical}, \\
    \text{human-review}, & \text{if } \tau_V > V_{t+1} \ge V_\mathrm{critical}, \\
    \text{none}, & \text{otherwise}.
  \end{cases}
\]
In addition, non-convergence or critical degradation is logged with a prime-specific rationale (e.g.\ $p=2 \times p=5$ interaction failure).

\subsection{Violation and Error Functionals}

On each batch of $N$ decisions, define:
\[
  \text{violations} = \frac{1}{N} \sum_{i=1}^N \mathbf{1}[\text{decision}_i \text{ is violation}],
\]
\[
  \text{interventions} = \frac{1}{N} \sum_{i=1}^N \mathbf{1}[\text{decision}_i \text{ requires human or governance override}].
\]

The error functional $E$ is taken as the violation rate (or a weighted variant) over a fixed horizon:
\[
  E = \mathbb{E}[\text{violations}] \approx \bar{v},
\]
and the reliability constraint is $E < \epsilon$.

The safety--efficiency optimization problem can be written as:
\[
  \min_{\Theta} \bigl(\text{violations}(\Theta), \text{interventions}(\Theta)\bigr)
  \quad \text{subject to} \quad E(\Theta) < \epsilon,
\]
with $\Theta$ including thresholds and update magnitudes.
This is the mathematical counterpart of tuning PMMA from Phase I (safety-max) to Phase II (frontier).

\section{Causal Validity Score}

\subsection{Model-Based Negative Log-Likelihood Gap}

To quantify causal validity $V_t$, PMMA uses the negative log-likelihood gap between a ground-truth causal model and a user-specified (possibly mis-specified) model.

Let $Y \in \{0,1\}$ be the outcome, $\bm{X}$ be observed predictors, and $C$ be a confounder included in the ground-truth model but potentially omitted by the user.
Let:
\begin{itemize}
  \item $p_\mathrm{gt}(Y \mid \bm{X}, C)$ be the ground-truth conditional probability.
  \item $p_\mathrm{user}(Y \mid \bm{X})$ be the user's model (e.g.\ omitting $C$).
\end{itemize}

Given a sample $\{(\bm{x}_i, c_i, y_i)\}_{i=1}^N$, define average negative log-likelihoods:
\[
  \mathrm{NLL}_\mathrm{gt} = -\frac{1}{N} \sum_{i=1}^N \bigl[y_i \log p_\mathrm{gt}(y_i \mid \bm{x}_i, c_i)
  + (1-y_i) \log(1-p_\mathrm{gt}(y_i \mid \bm{x}_i, c_i))\bigr],
\]
\[
  \mathrm{NLL}_\mathrm{user} = -\frac{1}{N} \sum_{i=1}^N \bigl[y_i \log p_\mathrm{user}(y_i \mid \bm{x}_i)
  + (1-y_i) \log(1-p_\mathrm{user}(y_i \mid \bm{x}_i))\bigr].
\]

The gap
\[
  \Delta = \max\{0, \mathrm{NLL}_\mathrm{user} - \mathrm{NLL}_\mathrm{gt}\}
\]
measures how much worse the user graph performs relative to the ground truth.
We map this gap to a validity score using an exponential decay:
\[
  V = \exp(-\lambda \Delta),
\]
for some scale $\lambda > 0$ (e.g.\ $\lambda = 5$).
Then $V \in (0,1]$, with $V \approx 1$ when the user graph is close to the ground truth and lower when confounders or edges are mis-specified.

\subsection{Prime $p=5$ and No-Code Risk}

In PMMA's prime-indexed interpretation, $p=5$ governs the tension between causal validity and the risks of no-code model specification.
A high mis-specification probability (user omits confounder $C$ often) leads to systematically higher $\Delta$ and lower $V$, triggering more human review and reducing autonomy.

\section{Example Code Snippets}

\subsection{Causal Validity Score}

\begin{verbatim}
import numpy as np

def nll_bernoulli(y, p):
    p = np.clip(p, 1e-6, 1 - 1e-6)
    return -np.mean(y * np.log(p) + (1 - y) * np.log(1 - p))

def validity_from_nll_gap(nll_user: float,
                          nll_ground_truth: float,
                          scale: float = 5.0) -> float:
    gap = max(0.0, nll_user - nll_ground_truth)
    return float(np.exp(-scale * gap))
\end{verbatim}

\subsection{Exception Router}

\begin{verbatim}
def route_exception(v5: float,
                    threshold: float,
                    converged: bool,
                    critical_threshold: float = 0.5):
    if not converged:
        return "human_on_exception", "non_convergence"
    if v5 < critical_threshold:
        return "human_on_exception", "critical_causal_degradation"
    if v5 < threshold:
        return "human_review", "causal_validity_gate"
    return "autonomous_ok", "none"
\end{verbatim}

\subsection{Budget Update}

\begin{verbatim}
def update_budget(current: float,
                  violation_pressure: float,
                  validity_pressure: float,
                  replenish: float = 0.0) -> float:
    next_budget = (current
                   - 0.10 * violation_pressure
                   - 0.05 * validity_pressure
                   + replenish)
    return max(0.0, min(1.0, next_budget))
\end{verbatim}

\subsection{Tuned PMMA Transition and Policy (Simplified)}

\begin{verbatim}
import numpy as np

def pmma_transition_tuned(state, violation_pressure, validity_pressure):
    autonomy, governance, budget = state

    base_step = 0.05
    extra_from_violation = 0.25 * violation_pressure
    extra_from_validity  = 0.20 * validity_pressure
    raw_delta = base_step * (extra_from_violation + extra_from_validity)
    violation_impact = min(0.08, raw_delta * (1.0 - governance))
    new_governance = np.clip(governance + violation_impact, 0, 1)

    autonomy_penalty = 0.20 * violation_pressure + 0.10 * validity_pressure
    autonomy_boost   = 0.05 * max(0.0, budget - 0.5)
    new_autonomy = np.clip(autonomy - autonomy_penalty + autonomy_boost, 0, 1)

    budget_decay = 0.10 * violation_pressure + 0.05 * validity_pressure
    budget_replenish = 0.02 if violation_pressure < 0.01 and validity_pressure == 0 else 0.0
    new_budget = np.clip(budget - budget_decay + budget_replenish, 0, 1)

    return np.array([new_autonomy, new_governance, new_budget])

def propose_autonomy_tuned(score_vec, v5, governance):
    high_score   = score_vec > 70
    good_causal  = v5 > 0.75
    governance_ok = governance < 0.60
    return high_score & good_causal & governance_ok

def pmma_policy_tuned(df, v5, base_autonomy_level,
                      base_governance_strictness,
                      delta=0.05, k_max=10, tau_v=0.60, v5_critical=0.50,
                      rng=None):
    if rng is None:
        rng = np.random.default_rng()

    state = np.array([base_autonomy_level, base_governance_strictness, 1.0])
    history = [state.copy()]
    score = np.clip(100 * df["p_true"].values, 0, 100)

    for step in range(k_max):
        autonomy, governance, budget = state

        proposed_auto = propose_autonomy_tuned(score, v5, governance)
        high_certainty = score > 90
        contradictions = proposed_auto & high_certainty & (governance >= 0.90)

        invalid_causal_action = proposed_auto & (v5 < tau_v)

        violations = (proposed_auto & (v5 < 0.70)) | contradictions
        intervention = (~proposed_auto) | (v5 < tau_v) | contradictions

        critical_escalation = (v5 < v5_critical) and violations.any()

        violation_pressure = violations.mean()
        validity_pressure = max(0.0, tau_v - v5)

        new_state = pmma_transition_tuned(state,
                                          violation_pressure,
                                          validity_pressure)
        dist = float(np.linalg.norm(new_state - state))
        history.append(new_state.copy())
        state = new_state

        if dist <= delta:
            return {
                "score": score,
                "autonomous_action": proposed_auto & ~(intervention),
                "intervention": intervention or critical_escalation,
                "contradictions": contradictions,
                "violations": violations or critical_escalation,
                "converged": True,
                "steps": step + 1,
                "state_history": np.array(history),
                "final_state": state,
                "failed_prime": None,
                "critical_escalation": critical_escalation,
            }

    return {
        "score": score,
        "autonomous_action": np.zeros(len(score), dtype=bool),
        "intervention": np.ones(len(score), dtype=bool),
        "contradictions": np.zeros(len(score), dtype=bool),
        "violations": np.zeros(len(score), dtype=bool),
        "converged": False,
        "steps": k_max,
        "state_history": np.array(history),
        "final_state": state,
        "failed_prime": "p=2 x p=5 non-convergence",
        "critical_escalation": np.array([v5 < v5_critical]),
    }
\end{verbatim}

\section{Safety--Efficiency Manifold and Phase I vs.\ Phase II}

The PMMA architecture admits distinct regimes depending on parameter choices:
\begin{itemize}
  \item \textbf{Phase I (Safety-Max):} conservative thresholds and aggressive governance updates yield low or zero baseline violations but very high intervention ratios and error budgets $E$ that may exceed the target $\epsilon$.
  \item \textbf{Phase II (Frontier-Seeking):} tuned thresholds and softer dynamics reduce violations below the target (e.g.\ $<5\%$) while lowering interventions to enterprise-acceptable levels (e.g.\ $<25\%$), maintaining $E < \epsilon$.
\end{itemize}

The same state space $M_t$ and recursion map $\Pi$ are used in both phases; only the configuration $\Theta$ changes.
This establishes PMMA not as a fixed solution but as a \emph{manifold of operating points} parameterized by explicit governance levers.

\section{Prior Art and Contribution}

Existing work on agentic AI governance, causal AI, and enterprise risk frameworks has:
\begin{itemize}
  \item Emphasized error budgets, liability tiers, and tiered human oversight (human-in-the-loop vs.\ human-on-the-loop).
  \item Highlighted the risks of causal mis-specification and black-box autonomy.
  \item Called for governance-by-design architectures where constraints are coded into orchestration frameworks.
\end{itemize}

However, those works generally do not:
\begin{itemize}
  \item Formalize these tensions as a prime-indexed multiplicity space with recursive stabilization.
  \item Define a causal validity gate as a continuous function of negative log-likelihood gap and embed it in a stateful agent recursion.
  \item Treat the system as a tunable safety--efficiency manifold explicitly parameterized by policy levers and evaluated via a simulation harness.
\end{itemize}

This defensive publication discloses:
\begin{enumerate}
  \item The PMMA state formalism and prime-indexed interpretation.
  \item The admissible recursion condition $\Lambda_t = \mathbf{1}[V_t \ge \tau_V]$ and non-convergence fallback.
  \item A concrete error budget functional $E$ and its use as a constraint in a bi-objective safety--efficiency optimization.
  \item A causal validity score $V$ derived from NLL gaps and used both as a gate and as a feedback signal in the recursion.
  \item Example code for budget updates, exception routing, and tuned PMMA policies.
\end{enumerate}

\section{Recommendations and Extensions}

\subsection{Fastest Path to Empirical Validation}

To validate PMMA in practice:
\begin{enumerate}
  \item Implement the simulation core for a single bounded domain (e.g.\ lead qualification).
  \item Run multi-seed experiments for baseline, Phase I, and Phase II policies under varying confounder omission rates.
  \item Measure violations per 100 decisions, intervention ratios, error budget $E$, convergence frequency, and exception routing statistics.
  \item Plot safety--efficiency frontiers and select operating points that match enterprise constraints.
\end{enumerate}

\subsection{Potential Extensions}

Future work may:
\begin{itemize}
  \item Add additional primes for new tensions (e.g.\ transparency vs.\ proprietary dependencies).
  \item Model multi-agent interactions where each agent has its own PMMA state and governance coupling terms.
  \item Integrate PMMA into real orchestration frameworks with runtime telemetry feeding back into the recursion.
\end{itemize}

\section*{Conclusion}

PMMA provides a governance-first, mathematically explicit architecture that instantiates the Phase Mirror philosophy: dissonances are not resolved but managed through mechanisms, budgets, and escalation policies.
By treating autonomy, governance, causal validity, and reliability as an evolving state subject to recursive stabilization, PMMA turns high-level governance desiderata into testable, tunable systems---and this document places those specific constructions into the public domain as prior art.

\appendix
\section*{Mathematical Appendix}

This appendix formalizes the Phase Mirror Multiplicity Agents (PMMA) architecture and records the operator norm bounds and basic stability results derived in the main body of the work.
The goals are:
\begin{enumerate}
  \item To make explicit the state space and operators underlying PMMA.
  \item To show under what conditions the recursive stabilization step is contractive in an appropriate norm.
  \item To derive bounds on autonomy, governance, and budget updates that ensure safety constraints such as $E < \epsilon$ can be enforced.
\end{enumerate}

\section{State Space and Operators}

\subsection{PMMA State Space}

Define the PMMA state at batch index $t$ as
\[
  M_t = (a_t, g_t, b_t, V_t, e_t) \in \mathcal{X} := [0,1]^4 \times \mathcal{E},
\]
where:
\begin{itemize}
  \item $a_t \in [0,1]$ is the autonomy level.
  \item $g_t \in [0,1]$ is the governance strictness level.
  \item $b_t \in [0,1]$ is the reliability (error) budget remaining.
  \item $V_t \in [0,1]$ is the causal validity score.
  \item $e_t \in \mathcal{E} := \{\mathrm{none},\mathrm{human\_review},\mathrm{human\_on\_exception}\}$ is the exception routing state.
\end{itemize}
For most of this appendix we restrict attention to the continuous part
\[
  x_t := (a_t, g_t, b_t, V_t) \in X := [0,1]^4,
\]
and treat $e_t$ as a piecewise-constant label updated by threshold conditions (Section~\ref{sec:exception-routing}).

\subsection{Policy Parameters and Admissible Recursion}

Let
\[
  \Theta
  = (\tau_V, V_{\mathrm{crit}}, \epsilon, \delta, k_{\max},
     \gamma_v, \gamma_u, r_{\max}, \alpha_a, \alpha_g, \beta_a)
\]
denote the collection of PMMA policy parameters, where:
\begin{itemize}
  \item $\tau_V \in (0,1)$ is the causal validity threshold for admissible recursion.
  \item $V_{\mathrm{crit}} \in (0,\tau_V)$ is the critical validity threshold for forced human-on-exception.
  \item $\epsilon \in (0,1)$ is the reliability budget target (e.g.\ $5\%$).
  \item $\delta > 0$ is the convergence tolerance.
  \item $k_{\max} \in \mathbb{N}$ is the maximum number of recursive steps per batch.
  \item $\gamma_v, \gamma_u > 0$ are budget decay rates for violation pressure and validity pressure.
  \item $r_{\max} \ge 0$ is the maximum budget replenishment per step under ideal behavior.
  \item $\alpha_a, \alpha_g$ encode autonomy and governance sensitivities to violations and validity pressure.
  \item $\beta_a$ encodes budget-dependent autonomy restoration.
\end{itemize}

Define the \emph{admissible recursion indicator}
\[
  \Lambda_t := \mathbf{1}[V_t \ge \tau_V].
\]
If $\Lambda_t = 0$, the system is not allowed to perform autonomous recursion; it must route decisions to human-review or human-on-exception.

\subsection{Violation and Validity Pressures}

On batch $t$, let
\[
  v_t := \frac{1}{N_t} \sum_{i=1}^{N_t} \mathbf{1}[\text{decision}_i \text{ is violation}],
\]
where $N_t$ is the batch size and ``violation'' is defined by policy (e.g.\ autonomous action with $V_t < 0.7$, or violation of governance rule at high certainty).
Let
\[
  u_t := \max\{0, \tau_V - V_t\}
\]
be the \emph{validity pressure}, i.e.\ how far below the recursion threshold the current validity lies.

\section{Core Operators}

\subsection{Transition Map on Continuous State}

We model one recursive step on the continuous state as the map
\[
  F: X \times \Theta \times [0,1]^2 \to X,
\]
\[
  x_{t+1} = F(x_t; \Theta, v_t, u_t).
\]
We specify $F$ component-wise:
\begin{align}
  g_{t+1}
    &= \mathrm{clip}_{[0,1]}\Bigl(
         g_t + \alpha_g(v_t, u_t)\,(1 - g_t)
       \Bigr),
       \label{eq:g-update}\\[4pt]
  a_{t+1}
    &= \mathrm{clip}_{[0,1]}\Bigl(
         a_t - \alpha_a(v_t, u_t) + \beta_a(b_t)
       \Bigr),
       \label{eq:a-update}\\[4pt]
  b_{t+1}
    &= \mathrm{clip}_{[0,1]}\Bigl(
         b_t - \gamma_v v_t - \gamma_u u_t + r_t
       \Bigr),
       \label{eq:b-update}\\[4pt]
  V_{t+1}
    &= V_t, \quad \text{(within a single recursion episode)}
       \label{eq:V-update-episode}
\end{align}
where $\mathrm{clip}_{[0,1]}$ denotes projection onto $[0,1]$, $r_t$ is the replenish term
\[
  r_t :=
  \begin{cases}
    r_{\max}, & \text{if } v_t < v_{\mathrm{tol}} \text{ and } u_t = 0,\\
    0, & \text{otherwise},
  \end{cases}
\]
and $v_{\mathrm{tol}}$ is a small tolerance for ``zero'' violation (e.g.\ $v_{\mathrm{tol}} = 0.01$).

For clarity, we may parameterize the sensitivities as:
\begin{align}
  \alpha_g(v_t,u_t)
    &= c_g^0 + c_g^v v_t + c_g^u u_t, \label{eq:alpha-g}\\
  \alpha_a(v_t,u_t)
    &= c_a^v v_t + c_a^u u_t, \label{eq:alpha-a}
\end{align}
with nonnegative coefficients $c_g^0, c_g^v, c_g^u, c_a^v, c_a^u$ chosen such that $\alpha_g,\alpha_a$ remain bounded on $[0,1]^2$.

The autonomy boost term $\beta_a$ is chosen to depend only on budget:
\[
  \beta_a(b_t) := c_b \max\{0, b_t - b_{\mathrm{mid}}\}
\]
for some $c_b \ge 0$ and $b_{\mathrm{mid}} \in (0,1)$.

\subsection{Budget Operator Norm Bounds}

We first bound the budget operator in the $\ell_\infty$ and $\ell_2$ norms.

\begin{lemma}[Budget Operator Lipschitz Bound]
\label{lem:budget-lip}
Fix $\gamma_v,\gamma_u, r_{\max}$ and define the pre-clipping budget update
\[
  \tilde{b}_{t+1}
    = b_t - \gamma_v v_t - \gamma_u u_t + r_t.
\]
Then, for any admissible $(b_t,v_t,u_t)$ and $(b_t',v_t',u_t')$ in $[0,1]^3$, we have
\[
  |\tilde{b}_{t+1} - \tilde{b}_{t+1}'|
  \le |b_t - b_t'|
     + \gamma_v |v_t - v_t'|
     + \gamma_u |u_t - u_t'|.
\]
Moreover, the clipped update $b_{t+1} = \mathrm{clip}_{[0,1]}(\tilde{b}_{t+1})$ satisfies the same bound, so the budget component of $F$ is Lipschitz with constant at most $1 + \gamma_v + \gamma_u$ with respect to $(b,v,u)$ in the $\ell_1$ norm.
\end{lemma}

\begin{proof}
By direct algebra,
\[
  \tilde{b}_{t+1} - \tilde{b}_{t+1}'
   = (b_t - b_t')
     - \gamma_v (v_t - v_t')
     - \gamma_u (u_t - u_t')
     + (r_t - r_t').
\]
Since $r_t,r_t' \in \{0,r_{\max}\}$, we have $|r_t - r_t'| \le r_{\max}$, but importantly $r_t$ is a function only of $(v_t,u_t)$ with bounded variation on a compact domain.
For Lipschitz control with respect to $(b,v,u)$, we can absorb $|r_t - r_t'|$ into $\gamma_v|v_t - v_t'| + \gamma_u |u_t - u_t'|$ by choosing $\gamma_v,\gamma_u$ large enough that any threshold discontinuity is dominated by their coefficients (or by restricting attention to regions where $r_t$ is constant).
Ignoring $r_t$ for the worst-case continuous regime,
\[
  |\tilde{b}_{t+1} - \tilde{b}_{t+1}'|
  \le |b_t - b_t'|
     + \gamma_v |v_t - v_t'|
     + \gamma_u |u_t - u_t'|.
\]
Clipping is a non-expansive map in every $\ell_p$ norm, hence the same inequality holds for $b_{t+1}$.
\end{proof}

\subsection{Full-State Operator Norm Bound}

We now treat the full continuous state operator $F$ and derive an operator norm bound in $\ell_2$.

\begin{theorem}[Global Lipschitz Bound for $F$]
\label{thm:F-lip}
Suppose $\alpha_g,\alpha_a,\beta_a$ are affine in $(v,u,b)$ as in \eqref{eq:alpha-g}--\eqref{eq:alpha-a} and the budget operator is as above.
Then there exists a constant $L_F > 0$ depending only on $\Theta$ such that for all $x_t,x_t' \in X$ and fixed $(v_t,u_t)$,
\[
  \| F(x_t;\Theta,v_t,u_t) - F(x_t';\Theta,v_t,u_t) \|_2
  \le L_F \, \|x_t - x_t'\|_2.
\]
In particular, we can take
\[
  L_F \le \sqrt{ L_a^2 + L_g^2 + L_b^2 + 1 },
\]
where $L_a,L_g,L_b$ are the component-wise Lipschitz constants for $a,g,b$ respectively.
\end{theorem}

\begin{proof}
We bound each component.
Write $x_t = (a_t,g_t,b_t,V_t)$ and similarly for $x_t'$.

\emph{Governance component.}
As a function of $(g_t,v_t,u_t)$, the pre-clipping update is
\[
  \tilde{g}_{t+1} = g_t + \alpha_g(v_t,u_t)(1-g_t).
\]
Holding $v_t,u_t$ fixed,
\[
  \frac{\partial \tilde{g}_{t+1}}{\partial g_t}
    = 1 - \alpha_g(v_t,u_t),
\]
so
\[
  |\tilde{g}_{t+1} - \tilde{g}_{t+1}'|
  \le |1 - \alpha_g(v_t,u_t)| \, |g_t - g_t'|.
\]
As $\alpha_g$ is bounded on $[0,1]^2$, say $|\alpha_g(v,u)| \le A_g$, we can take
\[
  L_g := \max\{ |1 - \alpha_g(v_t,u_t)| : (v_t,u_t)\in[0,1]^2 \}
  \le 1 + A_g.
\]
Clipping is non-expansive, so $|g_{t+1}-g_{t+1}'| \le L_g |g_t - g_t'|$.

\emph{Autonomy component.}
The pre-clipping update is
\[
  \tilde{a}_{t+1} = a_t - \alpha_a(v_t,u_t) + \beta_a(b_t).
\]
As a function of $(a_t,b_t)$ with $(v_t,u_t)$ fixed,
\[
  \frac{\partial \tilde{a}_{t+1}}{\partial a_t} = 1, \quad
  \frac{\partial \tilde{a}_{t+1}}{\partial b_t} = \beta_a'(b_t).
\]
Assuming $|\beta_a'(b)| \le B_a$ on $[0,1]$, we have
\[
  |\tilde{a}_{t+1} - \tilde{a}_{t+1}'|
  \le |a_t - a_t'|
     + B_a |b_t - b_t'|.
\]
Thus
\[
  |a_{t+1} - a_{t+1}'|
  \le L_a \sqrt{ (a_t-a_t')^2 + (b_t-b_t')^2 },
\]
with $L_a := \sqrt{1 + B_a^2}$.

\emph{Budget component.}
By Lemma~\ref{lem:budget-lip}, holding $(v_t,u_t)$ fixed,
\[
  |b_{t+1} - b_{t+1}'|
  \le L_b |b_t - b_t'|
\]
for some $L_b$ determined by $(\gamma_v,\gamma_u)$ (and uniform in $v_t,u_t$ on $[0,1]^2$).

\emph{Validity component.}
Within a recursion episode, $V_{t+1} = V_t$, so
\[
  |V_{t+1} - V_{t+1}'| = |V_t - V_t'|
\]
and the Lipschitz constant for this component is $1$.

Putting all components together and applying the triangle inequality and norm equivalence, we obtain
\[
  \| F(x_t;\Theta,v_t,u_t) - F(x_t';\Theta,v_t,u_t) \|_2
  \le \sqrt{L_a^2 + L_g^2 + L_b^2 + 1}\,\|x_t - x_t'\|_2.
\]
\end{proof}

\section{Contractivity and Recursive Stabilization}

\subsection{Single-Step Contractivity Conditions}

We now derive sufficient conditions for $F$ to be a contraction on $X$ under the $\ell_2$ norm.

\begin{corollary}[Contraction Sufficient Condition]
\label{cor:contraction}
If the parameter choices in $\Theta$ satisfy
\[
  \sqrt{L_a^2 + L_g^2 + L_b^2 + 1} < 1,
\]
then $F$ is a strict contraction on $X$ with respect to $\|\cdot\|_2$ for fixed $(v_t,u_t)$.
\end{corollary}

\begin{proof}
Immediate from Theorem~\ref{thm:F-lip}.
\end{proof}

In practice, it is difficult to make all components below unity in isolation, but note:
\begin{itemize}
  \item Governance updates can be made arbitrarily small by selecting $\alpha_g$ near zero when $g_t$ is large; this dampens $L_g$.
  \item Autonomy boosts can be limited by bounding $|\beta_a'|$ through a smooth, saturating function, so $B_a$ is small.
  \item Budget decay and replenishment can be tuned to keep $L_b$ close to $1$ or slightly below.
\end{itemize}
Together, this suggests a design procedure: choose $\Theta$ so that $F$ is non-expansive ($L_F \le 1$) in general, and contractive ($L_F < 1$) in the typical regime where violations and validity pressures are small.
This ensures that when the system is not under acute stress, the recursive update will stabilize quickly.

\subsection{Multi-Step Stability under Bounded Inputs}

Let $(v_t,u_t)$ be bounded inputs with $(v_t,u_t) \in [0,\bar{v}] \times [0,\bar{u}]$ for all $t$.
Assume $F$ is Lipschitz with constant $L_F \le 1$ and strictly contractive whenever $(v_t,u_t)$ lie in a ``good'' region $\mathcal{G}$ (low violations, high validity).

\begin{theorem}[Conditional Stabilization]
\label{thm:conditional-stab}
Suppose for all $t$ we have:
\begin{enumerate}
  \item $(v_t,u_t)\in\mathcal{G}$ implies $\|F(x_t;\Theta,v_t,u_t) - F(x_t';\Theta,v_t,u_t)\|_2 \le \rho \|x_t - x_t'\|_2$ for some $\rho \in (0,1)$.
  \item For $(v_t,u_t)\notin\mathcal{G}$, $F$ is non-expansive ($L_F \le 1$).
\end{enumerate}
Then for any initial $x_0,x_0'$, the distance $\|x_t - x_t'\|_2$ is non-increasing in $t$, and whenever the system spends a contiguous block of steps with $(v_t,u_t)\in\mathcal{G}$, the state contracts exponentially at rate at least $\rho$ on that block.
\end{theorem}

\begin{proof}
Non-expansiveness ensures $\|x_{t+1}-x_{t+1}'\|_2 \le \|x_t - x_t'\|_2$ for all $t$.
On any block $[s,s+T-1]$ where $(v_t,u_t)\in\mathcal{G}$ for all $t$ in that block, we have
\[
  \|x_{s+T} - x_{s+T}'\|_2 \le \rho^T \|x_s - x_s'\|_2.
\]
Therefore the distance is non-increasing globally and decays exponentially whenever the inputs remain in $\mathcal{G}$ for multiple steps.
\end{proof}

This theorem formalizes PMMA's \emph{conditional stabilization} interpretation: the system need not globally converge for all disturbance patterns, but when violation and validity pressures are within a good regime, the recursion tends to a stable fixed point; outside that regime, it is at least non-expansive and may shift to human-on-exception rather than attempting further autonomous recursion.

\section{Error Budget and Safety Constraint}

\subsection{Error Functional \texorpdfstring{$E$}{E}}

Let $v_t$ be the violation rate at batch $t$.
Over a horizon of $T$ batches, define the error functional
\[
  E_T := \frac{1}{T} \sum_{t=1}^T v_t.
\]
The reliability target is $E_T < \epsilon$.

\begin{lemma}[Error Budget under Bounded Violations]
\label{lem:error-budget}
Suppose $v_t \le v_{\max}$ for all $t$ and that $b_t$ is updated according to \eqref{eq:b-update} with $\gamma_v,\gamma_u,r_{\max}$ satisfying
\[
  \gamma_v v_{\max} + \gamma_u \max_t u_t \le \bar{\gamma}
  \quad\text{and}\quad
  r_{\max} \le \bar{r},
\]
with $\bar{\gamma},\bar{r}$ small.
Then there exists a choice of $(\bar{\gamma},\bar{r})$ such that whenever $E_T < \epsilon$,
the budget $b_t$ stays above a floor $b_{\min}(\epsilon)$ for all $t\le T$.
\end{lemma}

\begin{proof}
The worst-case budget change per step is
\[
  \Delta b_t = -\gamma_v v_t - \gamma_u u_t + r_t
  \ge -\gamma_v v_{\max} - \gamma_u \max_t u_t + 0
  \ge -\bar{\gamma}.
\]
Over $T$ steps, the total budget loss is at most $T\bar{\gamma}$.
Fix an initial $b_0$ and choose $\bar{\gamma}$ such that $b_0 - T\bar{\gamma} \ge b_{\min}(\epsilon)$, where $b_{\min}(\epsilon) > 0$ is an acceptable floor correlated with $\epsilon$.
Under $E_T < \epsilon$, we can ensure $v_{\max}$ and $\max_t u_t$ remain small (since low violations and adequate validity imply small pressures), making such a choice feasible.
\end{proof}

The practical implication is that the error budget dynamics can be tuned so that maintaining $E_T < \epsilon$ guarantees a non-degenerate budget, which in turn allows some autonomy to be preserved even under occasional violations.

\section{Exception Routing and Governability}
\label{sec:exception-routing}

We formalize exception routing as a function
\[
  R : [0,1]^2 \times \{0,1\} \to \mathcal{E}\times\mathcal{R},
\]
where $\mathcal{R}$ is a set of rationale codes.

Given $(V_{t+1},\Lambda_{t+1},\text{converged})$, define
\[
  (e_{t+1},r_{t+1}) = R(V_{t+1},\Lambda_{t+1},\mathrm{converged}),
\]
with:
\[
  R(V,\Lambda,\mathrm{converged}) =
  \begin{cases}
    (\mathrm{human\_on\_exception},\mathrm{non\_convergence}),
      & \text{if }\mathrm{converged} = 0,\\[4pt]
    (\mathrm{human\_on\_exception},\mathrm{critical\_causal\_degradation}),
      & \text{if } V < V_{\mathrm{crit}},\\[4pt]
    (\mathrm{human\_review},\mathrm{causal\_validity\_gate}),
      & \text{if } V_{\mathrm{crit}}\le V < \tau_V,\\[4pt]
    (\mathrm{none},\mathrm{none}),
      & \text{if } V \ge \tau_V \text{ and converged} = 1.
  \end{cases}
\]

\begin{definition}[Governability]
We say PMMA is \emph{governable} on horizon $T$ if for every batch $t\le T$, whenever $V_t < \tau_V$ or recursion fails to converge, we have $e_t \in \{\mathrm{human\_review},\mathrm{human\_on\_exception}\}$.
\end{definition}

\begin{proposition}[Sufficient Condition for Governability]
If exception routing $R$ is implemented as above and $\Lambda_t = \mathbf{1}[V_t \ge \tau_V]$, then PMMA is governable on every finite horizon $T$, regardless of the choice of $(v_t,u_t)$.
\end{proposition}

\begin{proof}
By construction, if $V_t < \tau_V$ then either $V_t < V_{\mathrm{crit}}$ or $V_{\mathrm{crit}} \le V_t < \tau_V$, and the routing cases enforce $e_t\in\{\mathrm{human\_review},\mathrm{human\_on\_exception}\}$.
If recursion fails to converge, the first case applies, also forcing $e_t = \mathrm{human\_on\_exception}$.
Thus, in all unsafe validity or convergence conditions, decisions are routed away from fully autonomous execution.
\end{proof}

In institutional governance terms, this proposition is the formal expression of ``knows when it should not act autonomously.''

\section{Prime-Indexed Decomposition}

Finally, we sketch how the above operators can be decomposed into prime-indexed components.

Let $\Pi = \{2,5\}$ for the two-prime pilot:
\[
  p=2 \quad\text{(autonomy vs.\ governance)},
  \qquad
  p=5 \quad\text{(causal validity vs.\ mis-specification)}.
\]
Define per-prime operators:
\begin{itemize}
  \item $A_2$ and $G_2$ acting on $(a,g)$ with parameters $(\alpha_a,\alpha_g)$.
  \item $C_5$ acting on $V$ with parameters $(\tau_V,V_{\mathrm{crit}})$ and data-dependent NLL gaps.
  \item $B$ acting on $b$ with parameters $(\gamma_v,\gamma_u,r_{\max})$ and pressures $(v,u)$.
\end{itemize}
Then one can write the overall transition as a composition
\[
  x_{t+1} = B \circ A_2 \circ G_2 \circ C_5 (x_t; \Theta,v_t,u_t),
\]
with each operator individually Lipschitz.
The global Lipschitz constant is bounded by the product of per-operator Lipschitz constants in the worst case, though the actual dynamics may be more benign.

This prime-indexed decomposition is useful conceptually because it allows per-tension analysis:
\begin{itemize}
  \item Tuning $A_2,G_2$ affects the $p=2$ autonomy--governance tradeoff, holding $C_5$ fixed.
  \item Tuning $C_5$ affects how aggressively the system responds to causal mis-specification, holding $A_2,G_2$ fixed.
\end{itemize}
In future work, additional primes can be added to capture new structural tensions while maintaining a modular analysis of operator norms and stability.

\section*{Summary}

This appendix has:
\begin{itemize}
  \item Defined the PMMA state space and transition operators.
  \item Derived Lipschitz bounds for the core transition map and provided sufficient conditions for non-expansiveness and contractivity.
  \item Related the error budget functional $E_T$ to the budget dynamics and showed how to keep budget above a floor under bounded violations.
  \item Formalized exception routing and proved a basic governability property.
  \item Sketched a prime-indexed operator decomposition for two-prime pilots.
\end{itemize}
These results provide the mathematical backbone for treating PMMA as a conditional stabilizer over a safety--efficiency manifold governed by explicit policy parameters rather than ad hoc heuristics.


\nocite{*}
\bibliographystyle{unsrt}  
\bibliography{references}  


\end{document}
